\documentclass[12pt,a4paper]{article}
\usepackage[utf8]{inputenc}
\usepackage[T2A]{fontenc}
\usepackage[russian]{babel}
\usepackage{amsmath,amssymb,amsthm}
\usepackage{booktabs,array}
\usepackage{geometry}
\usepackage{microtype}
\geometry{margin=2.3cm}

\newtheorem{proposition}{Предложение}
\newtheorem{nogocriterion}{Критерий провала}

\title{Parent-action gate для семейного смешивания:\\
коммутатор, projective flux и нечётномерное препятствие}
\author{}
\date{5 августа 2026 г.}

\begin{document}
\maketitle

\section{Что обязана дать новая вставка}

Факторные трансляции $T_{RP^3}$ и $T_{S^1}$ различают три семейных характера,
но коммутируют и потому дают общий собственный базис для up- и down-секторов.
Новый оператор должен одновременно:
\begin{enumerate}
 \item не коммутировать с факторной алгеброй;
 \item иметь фиксированную нормировку;
 \item следовать из parent action, а не из CKM-данных;
 \item порождать более одного mixing-направления.
\end{enumerate}

\section{Канонический commutator}

Из уже объявленных геометрических операторов можно построить эрмитов генератор
\begin{equation}
 K=\frac{i}{2}[T_{RP^3},S_{shear}].
\end{equation}
Его спектр равен $(-1,0,1)$, ранг равен двум, а нормировка не содержит
непрерывного параметра. Это первый прямой кандидат на некоммутирующую вставку.

Однако симметрийный тест отрицателен. Для группы факторных трансляций
\begin{equation}
 \frac14\sum_{g\in\mathbb Z_2^2}gKg^{-1}=0,
\end{equation}
и полное $AGL(2,2)$-усреднение также равно нулю. Следовательно, $K$ является
spurion-направлением нарушения симметрии, а не членом инвариантного действия.

Попытка управлять им через $SU(5)$-голономию,
\begin{equation}
 K\otimes P_{SU(5)},
 \qquad P_{SU(5)}=\operatorname{diag}(1,1,1,-1,-1),
\end{equation}
законна только после редукции $SU(5)$ к $S(U(3)\times U(2))$. Симметрия
разрешает такой член, но не фиксирует его коэффициент.

\begin{proposition}[Trace не равен action-selector]
Уникальность нормированного следа на $M_{60}(\mathbb C)$ фиксирует нормировку
уже записанного члена, но не относительные коэффициенты разных инвариантов.
На физическом триплете два неэквивалентных $SU(5)$-представления
$\mathbf{10}$ и $\overline{\mathbf5}$ уже допускают два независимых
квадратичных блока, то есть один относительный коэффициент после удаления
общего масштаба. Up- и down-Yukawa-инварианты также независимы.
\end{proposition}

Кроме того, один $K$ действует лишь в одной двумерной семейной плоскости и
оставляет одно направление точно изолированным. Он может породить один угол,
но не полную CKM-матрицу и не физическую CP-фазу.

\section{Нетривиальный $\mathbb Z_2$ projective flux}

Более сильная возможность возникает, если факторные трансляции реализованы
проективно:
\begin{equation}
 U^2=V^2=I,
 \qquad UV=-VU.
\end{equation}
Знак $-1$ является дискретным cocycle и даёт точную некоммутативность без
непрерывной настройки. На полном четырёхсостоянийном меню такая реализация
существует как пара signed permutation matrices.

Но здесь возникает простое нечётномерное препятствие. Если $U,V$ обратимы и
$UV=-VU$ в размерности $d$, то
\begin{equation}
 \det(UV)=(-1)^d\det(VU).
\end{equation}
При нечётном $d$ это противоречит $\det(UV)=\det(VU)$. Поэтому нетривиальный
projective flux не имеет трёхмерного представления.

В явной четырёхмерной реализации одна магнитная трансляция смешивает uniform-
синглет с семейным триплетом с нормой leakage, равной $1$. Сжатие операторов
обратно на триплет разрушает projective relation и снова даёт редуцируемую
алгебру с коммутантом размерности $2$, то есть структуру $1+2$.

\begin{nogocriterion}[Несовместимость flux и механизма $1+3$]
Нетривиальный $\mathbb Z_2$-cocycle является честным источником
коэффициент-свободной некоммутативности, но его представления чётномерны. Он
не может одновременно сохранять uniform reference state и действовать на
трёх поколениях. Искусственное утяжеление четвёртого состояния вводит новый
масштаб и scale-dependent Schur complement.
\end{nogocriterion}

\section{Новая развилка}

Текущий parent action не выводит CKM-вставку. Теперь остаются две существенно
разные программы:
\begin{enumerate}
 \item сохранить механизм $C^4=1+3$ и искать более богатый incidence/Dirac-
 оператор, который фиксирует минимум два spurion-направления;
 \item принять projective flux как первичный и заново вывести три наблюдаемых
 поколения из чётномерного носителя, отказавшись от прежнего $1+3$ аргумента.
\end{enumerate}
Смешивать эти варианты постфактум запрещено: они используют несовместимые
механизмы generation count.

Машинные протоколы сохранены в
\texttt{s2t\_parent\_noncommuting\_family\_insertion\_audit.py},
\texttt{s2t\_parent\_noncommuting\_family\_insertion\_results.json},
\texttt{s2t\_projective\_family\_flux\_audit.py} и
\texttt{s2t\_projective\_family\_flux\_results.json}.

\end{document}